\section{Introduction}
\label{sec:introduccion-ch}

L'hypothèse du continu demande s'il existe un ensemble de nombres réels dont la cardinalité soit strictement supérieure à celle des entiers naturels et inférieure à celle de la droite complète. Ce travail obtient une réponse négative à cette existence, et donc affirmative à l'hypothèse, dans une clôture généalogique définie au moyen des opérations de l'holographie modulaire triadique (HMT). Sa contribution relie la construction de la droite par prolongements compatibles à la comparaison cardinale de tous ses sous-ensembles internes.

Le noyau formel HMT réunit trois objets qui sont développés à partir de leurs
définitions. All Possible Paths (APP) munit le support
\((\mathbb Z/9\mathbb Z)^2\) des évaluations somme et produit, avec leurs
résidus et leurs quotients. Les translations par \((\pm1,0)\) et \((0,\pm1)\)
déterminent son graphe de Cayley ; les cellules carrées et l’identification des
bords opposés lui donnent une réalisation toroïdale bidimensionnelle. TRIT désigne
la signature ternaire orientée \(\{-1,0,+1\}\), qui distingue les
régimes des transitions. Le Topological Phase Kernel (TPK), ou noyau
topologique de phase, compose sélection, transport et mémoire pour prolonger
les états. Ces fonctions sont distinctes et agissent conjointement : le support
situe les chemins, les évaluations déterminent leur contenu arithmétique et la
dynamique conserve l’histoire de chaque transition.

L’argument suit cette composition depuis les opérations finies jusqu’au
domaine dans lequel s’effectue la comparaison cardinale :
\[
 \mathrm{APP}\longrightarrow\mathrm{TRIT}\longrightarrow\mathrm{TPK}
 \longrightarrow X_\infty^{\mathrm{enr}}
 \longrightarrow\mathcal C_{\mathrm{cont}}^{\mathrm{disc}}
 \longrightarrow m_\infty
 \longrightarrow\mathfrak G_{\mathrm{HMT}}(h,m_\infty).
\]
La composition produit des états enrichis, des prolongements compatibles,
des histoires complètes et, seulement ensuite, une réalisation archimédienne. La droite
réelle est une sortie de la construction, non son univers de départ ; la question
cardinale se décide sur le domaine intégral produit par cette généalogie.

La question centrale est de savoir si la droite obtenue par projection archimédienne admet, à l'intérieur de sa clôture généalogique complète, un cardinal strictement intermédiaire entre \(\aleph_0\) et le cardinal de la droite elle-même. Le théorème est formulé dans \(\mathfrak G_{\mathrm{HMT}}(h,m_\infty)\), la clôture engendrée par le système complet. Dans ce domaine, il quantifie sur \(\mathcal P^{\mathfrak G}(\mathbb R_{\mathrm{HMT}}^{\mathfrak G})\) et établit la dichotomie cardinale complète.

Le théorème~\ref{thm:decision-continuo-hmt}, démontré dans la
section~\ref{sec:cierre-cardinal-nativo}, établit le résultat principal.
Dans les comparaisons cardinales, la droite et les ensembles des parties sont
interprétés dans la clôture généalogique
\(\mathfrak G=\mathfrak G_{\mathrm{HMT}}(h,m_\infty)\). Ce domaine et ses
quantificateurs sont construits explicitement avant d’effectuer la comparaison.
Le résultat est le suivant. Pour toute histoire réalisée \(m_\infty\)
produite par le système complet associé à
\[
 \mathrm{APP}\longrightarrow\mathrm{TRIT}\longrightarrow\mathrm{TPK}
 \longrightarrow X_\infty^{\mathrm{enr}}
 \longrightarrow\mathcal C_{\mathrm{cont}}^{\mathrm{disc}},
\]
dont les restrictions sont compatibles et dont les prolongements sont non vides,
exhaustifs et survivants à tout horizon, l’extension généalogique
\(\mathfrak G_{\mathrm{HMT}}(h,m_\infty)\) engendre la droite intégrale,
quantifie sur tout son ensemble des parties et vérifie
\[
 \mathfrak G_{\mathrm{HMT}}(h,m_\infty)
 \models 2^{\aleph_0}=\aleph_1.
\]
En particulier, pour toute \(A\in\mathcal P^{\mathfrak G}(\mathbb R_{\mathrm{HMT}}^{\mathfrak G})\),
\[
 |A|\leq\aleph_0
 \qquad\text{ou}\qquad
 |A|=|\mathbb R_{\mathrm{HMT}}^{\mathfrak G}|.
\]
Par conséquent, HMT démontre affirmativement l'hypothèse du continu pour le continu intégral qu'elle définit et engendre. L'énoncé formel et sa preuve sont donnés une seule fois dans le théorème principal.

La différence essentielle par rapport à une présentation purement extensionnelle
réside dans l’identité de l’objet mathématique. Une coordonnée réelle isolée
conserve une valeur ; un état HMT conserve en outre les opérations qui l’ont
produit, le feuillet arithmétique, l’orientation tritique, les restes et
quotients traversés, la retenue, la frontière, la route et la mémoire accumulée.
La projection archimédienne ne définit ni n’épuise cette structure : elle détermine l’une de
ses coordonnées. La seconde projection conserve l’incidence de frontière qui,
dans des développements descendants, atteint Paley--Witt--Golay--Leech,
\(V^\natural\) et le Monstre. Les deux réalisations proviennent du même état et
aucune n’est introduite pour sélectionner rétrospectivement l’autre.

\subsection{Construction du continu par prolongements compatibles}

APP dispose sur un même support les feuillets somme et produit et conserve la
différence entre leurs évaluations additive et multiplicative. TRIT incorpore l’orientation et le régime. TPK
sélectionne, transporte, mémorise et prolonge les transitions admissibles. La
composition ne s’achève pas dans un état fini : elle définit un système de restrictions
\[
 \rho_{N+1,N}:X_{N+1}\longrightarrow X_N
\]
et des extensions non vides
\[
 \operatorname{Ext}_N(x)\subseteq\rho_{N+1,N}^{-1}(x).
\]
La compatibilité à toute profondeur produit la limite projective
\[
 X_\infty^{\mathrm{enr}}
 =\varprojlim_N(X_N,\rho_{N+1,N}).
\]
Le quantificateur est \(\forall N\) ; aucune exécution à mille, dix mille ou un million de chiffres ne constitue la borne du théorème.

Chaque bloc successeur subdivise un cylindre en \(729=3^6\) sous-cylindres de raffinement disjoints, exhaustifs et de diamètre décroissant. Ce nombre décrit la ramification locale du raffinement ; il n'est pas identifié au cardinal du continu. Le ruban complet de blocs détermine une histoire et sa projection arithmétique détermine un unique point réel. Réciproquement, le développement canonique de chaque réel reconstruit un ruban et la propriété de couverture le relève en une histoire HMT compatible. Ainsi, la droite est l'image surjective de l'arbre construit, jamais un ensemble extérieur utilisé pour choisir ses branches.

La mesure intervient comme couplage. Sur le domaine considéré, on
impose d’abord la condition de séparation des fibres : le lecteur enrichi
doit distinguer deux états différents qui produisent la même valeur sous l’évaluation archimédienne. Sous cette
condition, formalisée dans la
proposition~\ref{prop:ch-recuperabilidad-fibra}, l’instrument distingue une
continuation compatible, ajoute au registre généalogique le résultat et la
rétroaction et restreint les histoires futures à celles qui prolongent le nouvel
état. L’histoire réalisée \(m_\infty\) se forme au moyen de ces
mises à jour. La transformation inverse reconstruit alors l’intérieur
à partir de la mémoire transférée à la frontière, de sorte que la mise à jour
n’équivaut pas à une perte d’information. La
proposition~\ref{prop:ch-naturalidad-medida} montre quand les changements de
représentation conservent également la mesure des histoires.

\subsection{Clôture sémantique généalogique}

Soit \(h\) le code effectif des graphes d’APP, de TRIT, du TPK, des
prolongements, des restrictions et des deux projections. Le couple
\((h,m_\infty)\) est encodé primitivement au moyen d’un unique réel
\(u_{h,m}=h\oplus m_\infty\). On définit
\[
 \begin{aligned}
 \mathfrak G_0(h,m_\infty)
   &:=\operatorname{TC}\{\omega,u_{h,m},h,m_\infty\},\\
 \mathfrak G_{\beta+1}(h,m_\infty)
   &:=\operatorname{Def}_{\Sigma_{\mathrm{HMT}}}
        \bigl(\mathfrak G_\beta(h,m_\infty)\bigr),\\
 \mathfrak G_\lambda(h,m_\infty)
   &:=\bigcup_{\beta<\lambda}\mathfrak G_\beta(h,m_\infty),\\
 \mathfrak G_{\mathrm{HMT}}(h,m_\infty)
   &:=\bigcup_{\beta\in\operatorname{Ord}}
        \mathfrak G_\beta(h,m_\infty).
 \end{aligned}
\]
Cette clôture est formulée au moyen de la définissabilité relative et de la condensation au
sens précis des hiérarchies constructibles~\cite{Jech2003}. L’hypothèse
du continu n’est pas introduite parmi les conditions qui définissent la clôture ; elle
s’obtient au moyen de la démonstration développée dans l’article. On n’introduit pas non plus
comme données une énumération de la droite ni une bijection cible. Son
opération successeur forme les sous-ensembles définissables au moyen de formules finies
du langage HMT et de paramètres déjà produits.
De manière équivalente, tout objet admis possède un arbre de dérivation et tout
arbre admissible détermine un objet.

La hiérarchie résultante est transitive, contient ses ordinaux, satisfait les axiomes employés dans la comparaison cardinale et possède un bon ordre global déterminé par le premier niveau de naissance et le plus petit nom de dérivation. La notation
\[
 \mathfrak G_{\mathrm{HMT}}(h,m_\infty)
 =L[u_{h,m}]=L[h,m_\infty]
\]
s'obtient ensuite comme théorème de représentation. Elle ne constitue pas le moteur d'APP, de TRIT ou de TPK, ne sélectionne pas \(m_\infty\) et ne remplace pas la construction matérielle du continu.

\subsection{Comparaison cardinale entre la droite HMT et le premier ordinal non dénombrable de la clôture}

La preuve cardinale se divise en deux directions indépendantes de l’égalité
recherchée. D’abord, la condensation généalogique démontre que chaque réel engendré
apparaît avant \(\omega_1^{\mathfrak G}\). Pour maintenir la même
convention que dans la preuve, soient \(\iota(r)\) le code de la coupure rationnelle
de \(r\), \(\beta(r)\) le plus petit ordinal tel que
\(\iota(r)\in\mathfrak G_{\beta(r)+1}\), et \(n(r)\) le plus petit indice
de cette coupure dans l’énumération \(e_{\beta(r)+1}\). On définit
\[
 j_{\mathrm{HMT}}(r)=\omega\beta(r)+n(r),
 \qquad
 j_{\mathrm{HMT}}:
 \mathbb R_{\mathrm{HMT}}^{\mathfrak G}
 \hookrightarrow\omega_1^{\mathfrak G}.
\]

Dans la direction opposée, chaque ordinal \(\beta<\omega_1^{\mathfrak G}\) possède un code réel \(w_\beta\) d'un bon ordre dénombrable de type \(\beta\). Les codes sont choisis au moyen du bon ordre généalogique déjà construit. Le plongement ternaire
\[
 E(w)=\sum_{n\geq0}\frac{2\,\mathbf 1_w(n)}{3^{2n+2}}
\]
est injective, et le relèvement uniforme des cylindres prouve que
\(E(w_\beta)\) appartient à la droite obtenue par la projection archimédienne HMT. Par conséquent,
\[
 k_{\mathrm{HMT}}:
 \omega_1^{\mathfrak G}
 \hookrightarrow\mathbb R_{\mathrm{HMT}}^{\mathfrak G}.
\]
Cantor--Bernstein clôt l’égalité cardinale. L’identification démontrée
\[
 \mathbb R_{\mathrm{HMT}}^{\mathfrak G}
 \approx\mathcal P^{\mathfrak G}(\omega)
\]
la transforme exactement en \(2^{\aleph_0}=\aleph_1\). Ici, \(A\approx B\) signifie équipotence ; il n'affirme pas un isomorphisme de corps entre une droite et un ensemble des parties.

\clearpage
\subsection{Structure logique de la démonstration}

La démonstration s’organise en quatre étapes logiques. D’abord, on construit le
domaine des histoires, sa limite inverse et la couverture exhaustive des cylindres.
Ensuite, on forme la clôture sémantique et on applique la condensation à ses réels.
Troisièmement, on définit, sans utiliser l’égalité recherchée, les deux injections entre la
droite HMT et \(\omega_1^{\mathfrak G}\). Quatrièmement, Cantor--Bernstein conclut à
l’égalité et l’argument de cofinalité établit la dichotomie pour chaque
sous-ensemble interne. Dans cette séquence, \(729\) décrit la ramification locale ;
la couverture relie histoires et valeurs ; la condensation contrôle le niveau
d’apparition des réels internes. Les constantes sont des sorties
corrélées du même état : leurs chiffres conventionnels ne sélectionnent pas la
comparaison cardinale, tandis que leurs régions, routes et mémoires demeurent dans
l’objet enrichi sur lequel la droite est construite.
L’identification ultérieure à \(L[h,m_\infty]\) représente la clôture déjà
construite.

L'ordre d'exposition reproduit d'abord le noyau formel commun de la série, le prolongement, la construction conjointe du continu et la génération coinductive ; ensuite, la clôture sémantique intégrale et la preuve cardinale. Ce n'est qu'alors qu'il développe le noyau réduit de phase et de mémoire, le sceau dodécaphasique entier \(K\), la voie dodécaphasique complète de \(\alpha_{\mathrm{HMT}}\), le témoin analytique fini d'ordre neuf, la représentation analytique corrélée avec récurrence, convergence et cylindres communs, et le corridor exceptionnel. Ces coordonnées partagent un état d'origine, mais ont des codomaines et des fonctions différents ; ensemble, elles exhibent la double réalisation complète et conservent les chemins d'évaluation et de projection, l'incidence, la frontière et la mémoire. Aucune valeur cible ne sélectionne les injections cardinales. La comparaison avec Cantor, Gödel et Cohen est introduite au terme de la construction~\cite{Cantor1878,Hilbert1902,Godel1938,Cohen1963,Cohen1964} et ne gouverne pas la définition de l'objet HMT.

Chaque flèche utilisée par la clôture cardinale est définie dans le corps par son
domaine, son codomaine, son action et sa compatibilité avec le raffinement : les trois boucles
orientées du retour nonadique, leur regroupement en blocs de six symboles,
l’extension de bloc de \(54\) étapes, la construction des cylindres et le
relèvement uniforme de leurs branches. La construction explicite des neuf
transitions fournit le relèvement uniforme utilisé dans la comparaison
cardinale. Les calculs reproductibles accompagnent cette démonstration, mais ne
remplacent aucun de ses arguments.

% BEGIN R27_FRAMING_INTRO_ES
La transition des germes aux histoires compatibles contient des
lectures finies qu’il convient de distinguer par leur information. Les signatures
additive et multiplicative déterminent les \(42\) blocs possibles d’une
fenêtre décimale, avec des germes explicites qui prouvent leur atteignabilité.
Les émissions de six fenêtres établissent un cycle de raccords entre
moitiés régionales. Le transport temporel conserve, en plus de ces
incidences, la phase, l’orientation et la mémoire. Une paire de registres dodécaphasiques
ayant le même histogramme et une corrélation différente démontre pourquoi
l’archive exige l’ordre des échantillons. Ces preuves précisent
l’antécédent fini du prolongement ; le recouvrement, la clôture sémantique
et les injections cardinales sont développés dans leurs domaines propres.
% END R27_FRAMING_INTRO_ES

% BEGIN R28_FRAMING_INTRO_ES
La partie fondationnelle finie comprend le multigraphe d’émissions, son
quotient de phase et cinq clôtures minimales à huit arêtes. Chaque clôture
couvre conjointement APP, propagation, auto-échelle et une région de clôture.
Les témoins orientés, la recherche finie et la parité des lecteurs
de mémoire précisent les données conservées avant les prolongements
et les lectures cardinales développés dans l’ouvrage.
% END R28_FRAMING_INTRO_ES
